\section{Higher-order corrections for marked generating functions}\label{sec:higher-order}

The estimate of Proposition~\ref{prop:prefactor} has relative error
$O_p(\log g/g)$. The logarithm comes from integrating the error
$O_p(1/(hAr))$ in~\eqref{eq:PF11} towards the critical endpoint, where
$Ar$ has a simple zero. This section states two refinements, still for a
fixed number $p$ of distinguished markings: an explicit first correction,
which also gives the leading term with relative error $O_p(g^{-1})$, and a
complete expansion whose truncation order may grow like $\sqrt g$. Both
results are proved in Appendix~\ref{app:marked}; the passage to a
prescribed cusp count is made in
Sections~\ref{extracting-fixed-cusp-counts} and~\ref{sec:sharp-extraction}.

\subsection{The first correction}\label{subsec:first-correction}
The mechanism is that of Section~\ref{sec:prefactor}, carried one order
further. Lemma~\ref{lem:second-profile} extends the degree expansion
\eqref{eq:PF5} of the normalized profiles by a term $e_p(\mathbf d)/h^2$,
uniformly for $|\mathbf d|\le h^{1/8}$. Inserted into the dilaton identity,
it adds a term $\gamma_p/h$ to the numerator of the differential
equation~\eqref{eq:PF8}; see~\eqref{eq:ST8}. The coefficient $\gamma_p$,
defined in~\eqref{eq:ST7}, is a finite expression in kernel moments
(see~\eqref{eq:ST11}), and it vanishes at the pure-$\psi$ endpoint
(see~\eqref{eq:ST12}). Its primitive against $1/(Ar)$ is therefore bounded,
although $Ar$ vanishes there. The integration constant is fixed at the
endpoint by the first correction to $C_g(2^{3g-3})$, which we derive from
Painlev\'e~I in Subsection~\ref{sec:endpoint-expansion}, together with
the first correction to an explicit negative-binomial coefficient.
Accordingly, define
\begin{equation}
 t_p=\frac{-2p^2+12p-11}{20},\qquad
 \mathcal L_p(Y)=t_p-\int_Y^{y_0}\frac{\gamma_p(v)}{A(v)r(v)}\,\mathrm{d}v.
 \label{eq:ST13}
\end{equation}

\begin{theorem}[Uniform marked first correction]\label{thm:first-correction}
For every fixed $p\ge1$, uniformly for $0\le Y<y_0$,
\begin{equation}
 \Phi_g^{(p)}(X(Y))=\frac{\mathcal J_{g,p}(Y)}\pi
 \left\{1+\frac{\mathcal L_p(Y)}h+
            O_p\!\left(\frac{\log(g+1)}{g^2}\right)\right\}.
 \label{eq:ST14}
\end{equation}
In particular the leading relative error is $O_p(g^{-1})$.
\end{theorem}

At $Y=0$ the coefficients of this expansion reduce to the first-order
formulas of Mirzakhani and Zograf~\cite[Theorems~1.4 and~3.2(i)]{MZ};
this check is recorded at the end of Subsection~\ref{app:second-degree}.

\subsection{All orders}\label{subsec:all-orders}
The same scheme continues to every order. The recursion is triangular:
each new degree coefficient of the profiles uses earlier cusp and genus
coefficients, and the dilaton and KdV equations then give the new cusp and
genus coefficients. The difficulty is to keep the constants under control
as the order grows.

\emph{Where the squared factorial comes from.} The first two profile coefficients are polynomials in the degrees of degrees two and four: $\delta_p$ in~\eqref{eq:PF5} is quadratic in $D=|\mathbf d|$, and $e_p$ in~\eqref{eq:ST3} is quartic. This pattern persists. Each new order is obtained by summing an adjacent difference along a coordinate path of length $D$, and each difference involves pair counts of size $D$ and one factor $1/h$. So the order-$j$ coefficient behaves like a polynomial of degree $2j$ in the degrees, of size about $(D+j)^{2j}$ up to a factor $B^j$. In the dilaton sum the degree $d$ is weighted by the geometrically decaying coefficients $b_d$, so the effective size of $D$ is of order $j$, and $(Mj)^{2j}$ is $(j!)^2$ up to a geometric factor. The bound $B^{K+1}((K+1)!)^2h^{-K-1}$ for the first omitted term is smallest at an order proportional to $\sqrt h$, and a sufficiently small fixed multiple of $\sqrt h$ gives an exponentially small remainder. This explains both the Gevrey-2 bounds and the truncation order in Theorem~\ref{thm:full-expansion}.

Appendix~\ref{app:marked} constructs, by the finite triangular rules
\eqref{eq:FE1}--\eqref{eq:FE6}, the formal coefficients of the profiles
and of the cusp and genus ratios, and from them functions $a_{p,j}(Y)$ and
\[
 \mathcal L_{p,j}(Y)=[t^j]\exp\!\Bigl(\sum_{i=1}^j a_{p,i}(Y)t^i\Bigr),
 \qquad \mathcal L_{p,1}=\mathcal L_p;
\]
see~\eqref{eq:FE11}. On one complex neighborhood of $[0,y_0]$ they are
holomorphic and satisfy
$\|a_{p,j}\|+\|\mathcal L_{p,j}\|\le C_p\mathfrak B_p^j(j!)^2$.

\begin{theorem}[Complete uniform marked expansion]\label{thm:full-expansion}
For each fixed $p\ge1$ there are $C_p,\mathfrak B_p,c_p,h_p>0$ such that,
for $h=2g-2+p\ge h_p$ and $0\le K\le c_p\sqrt h$,
\begin{equation}
 \sup_{0\le Y<y_0}\left|
 \frac{\pi\Phi_g^{(p)}(X(Y))}{\mathcal J_{g,p}(Y)}
       -\sum_{j=0}^K\frac{\mathcal L_{p,j}(Y)}{h^j}\right|
 \le\frac{C_p\mathfrak B_p^{K+1}((K+1)!)^2}{h^{K+1}},
 \qquad \mathcal L_{p,0}=1.
 \label{eq:FE12}
\end{equation}
The coefficients are specified by the finite rules and integrals of
Appendix~\ref{app:marked}, and have the Gevrey-2 norms stated above on
one complex neighborhood of $[0,y_0]$. The same real error holds for the
logarithm with coefficients $a_{p,j}$. A sufficiently small square-root
truncation has error $O_p(\mathrm{e}^{-a_p\sqrt h})$; no convergence of
the infinite series is required.
\end{theorem}

The proof has three steps. Proposition~\ref{prop:hierarchy} divides each
order-$j$ coefficient by $(D+p+Mj)^{2j}$. After this division the
coefficient rules become coefficientwise inequalities between power
series with nonnegative coefficients, and a system of analytic majorants
(Lemma~\ref{lem:majorant}) bounds the normalized coefficients and the
actual remainders, uniformly in the order. Next, at the pure-$\psi$
endpoint, the integration constants are computed from the
Itzykson--Zuber identity and the large-order behavior of Painlev\'e~I;
the Borel representation of Costin, Costin and Huang~\cite{CCH} gives
these constants with a Gevrey-1 remainder. Finally the logarithmic
derivative is integrated from a matching point at distance $h^{-K-2}$
from the endpoint, where the positive homogeneous structure of the
amplitudes controls the change from the endpoint value.

At $Y=0$ this is the established fixed-cusp expansion mechanism;
effective order-dependent versions appear in
Hide--Macera--Thomas~\cite[Theorem~4.2]{HMT}. The estimates here remain
uniform up to the pure-$\psi$ endpoint. Their dependence on the order is
what permits the growing-order truncation in the absolute-volume theorem.
